\subsection{Protocol execution}
\label{subsec:protocol-execution}
This section provides an overview of the study conduction, following the protocol designed in Section \ref{subsec:protocol-definition}. Specifically, we first provide details about
the selection process and quality assessment of the included papers, followed by details on the data extraction process.


\subsubsection{Studies Selection Process} 
The process followed to select primary studies is illustrated in Figure \ref{fig:selection-process}. The Figure depicts the sequence of executed steps and their outcomes in terms of retrieved, discarded, and finally selected papers. 

\begin{enumerate}
\item \textit{Querying of digital libraries}: In this initial step, the search strings (introduced in Section \ref{subsec:protocol-definition}) were used to query the four digital libraries selected for our SLR (ACM, IEEE Xplore, Scopus, and WoS). No restrictions were applied on the publication year during the search, which was conducted on January 31st, 2025, resulting in a total of $1,370$ records.

\item \textit{Removal of duplicate records}: From the initial $1,370$ retrieved records, \textbf{316} duplicate entries were identified and removed, leaving $1,054$ studies for  screening.

\item \textit{Application of Inclusion and Exclusion Criteria (IC/EC) — first step}: The remaining 1,054 records were split into two groups, and their titles, abstracts, and metadata were independently screened by two researchers applying all IC and EC defined, except for EC6 and IC2, which require full-text analysis. During this phase, $912$ studies were excluded by mutual agreement, resulting in $142$ studies progressing to the next step.

\item \textit{Application of IC and EC on full texts}: The $142$ papers were divided between the two researchers, who performed a full-text assessment applying all the IC/EC including EC6 and IC2. This step resulted in the exclusion of $66$ studies and the retention of $76$ papers, of which $16$ were classified as "doubtful" due to disagreement between reviewers.

\item \textit{Resolution of doubtful studies}: To resolve the doubtful cases, two additional researchers joined the initial reviewers for a joint discussion and re-assessment based on full-text analysis. After this collaborative review, $14$ studies were excluded, and $2$ were included, resulting in a final set of $62$ selected studies.
\end{enumerate}

%\paragraph{Quality Assessment} 
%Table \ref{tab:QC-evaluation} summarizes the results of applying the Quality Criteria to the 41 selected papers. All studies achieved an overall quality score of 3 or higher out of a maximum of 6, indicating a satisfactory level of methodological rigor and reporting quality across the included literature.

\subsubsection{Quality assessment} 
Each of the $62$ selected papers was independently evaluated by two researchers using the six QCs described in Section \ref{subsec:protocol-definition}. Papers that received a total score below 3 from both researchers were excluded, while those with substantially differing scores or borderline cases (scores between 3 and 4) were discussed collectively by all researchers. Following this process, $41$ studies were retained as the final set of high-quality primary studies included in this review. The complete list of selected papers along with their quality scores is presented in Table \ref{tab:QC-evaluation}.

\subsubsection{Data extraction}  

Given that our RQs (see Section \ref{subsec:protocol-definition}) span the distinct domains of AI and marketing, the review team was composed of experts in both fields. The team collaborated to extract relevant data from each of the 41 selected studies, utilizing the data extraction form presented in Table \ref{tab:DEF}. In cases of uncertainty, disagreements were resolved through discussion to reach a consensus.

The data extraction form and the collected data are publicly accessible online.\footnote{\url{https://docs.google.com/spreadsheets/d/1d6RD3cnYkc6AYOfw-6SbhhxwG-dBsKlVmFowj5EHDg8/edit?usp=sharing}}  
%\DD{}{Add a link to the spreadsheet with the extracted evidence}


% Please add the following required packages to your document preamble:
% \usepackage{graphicx}
% \usepackage[table,xcdraw]{xcolor}
% Beamer presentation requires \usepackage{colortbl} instead of \usepackage[table,xcdraw]{xcolor}
\begin{table}[H]
\centering
\caption{Quality criteria assessment and resulting selected studies.}
\label{tab:QC-evaluation}
\resizebox{\textwidth}{!}{%
\begin{tabular}{|l|l|c|c|c|c|c|c|c|} 
\hhline{~~-------|}
\multicolumn{1}{l}{} &  & \multicolumn{7}{c|}{\cellcolor{lightgray}\raisebox{0.4\height}{\rule{0pt}{4ex}\large\textbf{Quality Criteria (QC)}}} \\ \hline
\rowcolor{lightgray}
\multicolumn{1}{|c|}{\cellcolor{lightgray}\raisebox{0.4\height}{\rule{0pt}{4ex}\large\textbf{ID}}} &
  \raisebox{0.4\height}{\rule{0pt}{4ex}\large\textbf{Study Title}} &
  \multicolumn{1}{c|}{\cellcolor{lightgray}\raisebox{0.4\height}{\rule{0pt}{4ex}\textbf{1}}} &
  \multicolumn{1}{c|}{\cellcolor{lightgray}\raisebox{0.4\height}{\rule{0pt}{4ex}\textbf{2}}} &
  \multicolumn{1}{c|}{\cellcolor{lightgray}\raisebox{0.4\height}{\rule{0pt}{4ex}\textbf{3}}} &
  \multicolumn{1}{c|}{\cellcolor{lightgray}\raisebox{0.4\height}{\rule{0pt}{4ex}\textbf{4}}} &
  \multicolumn{1}{c|}{\cellcolor{lightgray}\raisebox{0.4\height}{\rule{0pt}{4ex}\textbf{5}}} &
  \multicolumn{1}{c|}{\cellcolor{lightgray}\raisebox{0.4\height}{\rule{0pt}{4ex}\textbf{6}}} &
  \multicolumn{1}{c|}{\cellcolor{lightgray}\raisebox{0.4\height}{\rule{0pt}{4ex}\textbf{\parbox[c]{1cm}{\centering Total\\Score}}}} \\
  \hline
\textbf{S1}                                                                       & Customer data prediction and analysis in e-commerce using machine learning \cite{S1-ID3-Rahib}                                                                                                                                                          & 0.5        & 1          & 0.5        & 0          & 0.5        & 1          & \textbf{3.5}   \\ 
\hline
\rowcolor{lightgray} \textbf{S2}                                     & E-commerce big data processing based on an improved RBF model \cite{S2-ID16-Lu}                                                                                                                                                                       & 1          & 0.5        & 0.5        & 0          & 1          & 1          & \textbf{4}     \\ 
\hline
\textbf{S3}                                                                       & \begin{tabular}[c]{@{}l@{}}Comparative SHAP Analysis on SVM and K-NN: Impacts of Hyperparameter Tuning on \\ Model Explainability \cite{S3-ID223-Beckhauser}\end{tabular}                                                                                       & 0          & 1          & 1          & 0.5        & 1          & 1          & \textbf{4.5}   \\ 
\hline
\rowcolor{lightgray} \textbf{S4}                                     & \begin{tabular}[c]{@{}>{\cellcolor{lightgray}}l@{}}Analyzing Customer Churn Behavior Using Data mining Approach: Hybrid Support \\ Vector Machine and Logistic Regression in Retail Chain \cite{S4-ID40-Barzegar}\end{tabular}                  & 0.5        & 1          & 0.5        & 0          & 1          & 1          & \textbf{4}     \\ 
\hline
\textbf{S5}                                                                       & \begin{tabular}[c]{@{}l@{}}TSUNAMI - an explainable PPM approach for customer churn prediction in evolving \\ retail data environments \cite{S5-ID50-Pasquadibisceglie}\end{tabular}                                                                                  & 0.5        & 0.5        & 1          & 1          & 1          & 1          & \textbf{5}     \\ 
\hline
\rowcolor{lightgray} \textbf{S6}                                     & \begin{tabular}[c]{@{}>{\cellcolor{lightgray}}l@{}}Predicting customer churn in the fast-moving consumer goods segment of the retail \\ industry using deep learning     \cite{S6-ID70-Mahdi}\end{tabular}                                       & 0          & 1          & 0.5        & 0.5        & 0.5        & 1          & \textbf{3.5}   \\ 
\hline
\textbf{S7}                                                                       & \begin{tabular}[c]{@{}l@{}}Explainable machine learning models applied to predicting customer churn for \\e-commerce \cite{S7-ID79-Boukrouh}\end{tabular}                                                                                                   & 0.5        & 0.5        & 1          & 0          & 0.5        & 0.5        & \textbf{3}   \\ 
\hline
\rowcolor{lightgray} \textbf{S8}                                     & Prediction of Customer Behavior Changing via a Hybrid Approach                                                                                                                                                                       & 1          & 0.5        & 1          & 0.5        & 1          & 1          & \textbf{5}   \\ 
\hline
\textbf{S9}                                                                       & \begin{tabular}[c]{@{}l@{}}Vision Transformer Approach to Customer Churn Prediction Radar Chart Image \\ Classification for Non-subscription Based E-commerce \cite{S9-ID95-Coolwijk}\end{tabular}                                                           & 0          & 1          & 0.5        & 1          & 0.5        & 1          & \textbf{4}   \\ 
\hline
\rowcolor{lightgray} \textbf{S10}                                    & \begin{tabular}[c]{@{}>{\cellcolor{lightgray}}l@{}}Development of fading channel patch based convolutional neural network models for \\ customer churn prediction \cite{S10-ID98-Seema}\end{tabular}                                          & 0.5        & 1          & 1          & 0.5        & 1          & 1          & \textbf{5}   \\ 
\hline
\textbf{S11}                                                                      & \begin{tabular}[c]{@{}l@{}}Deciphering Digital Social Dynamics: A Comparative Study of Logistic Regression and \\ Random Forest in Predicting E-Commerce Customer Behavior \cite{S11-D105-Sunarya}\end{tabular}                                              & 0.5        & 1          & 0.5        & 0.5        & 0.5        & 1          & \textbf{4}   \\ 
\hline
\rowcolor{lightgray} \textbf{S12}                                    & A Predictive Model for E-Commerce Customer Churn under an Intelligent Algorithm \cite{S12-ID198-Tuo}                                                                                                                                                     & 0          & 0.5        & 0.5        & 1          & 0.5        & 0.5        & \textbf{3}     \\ 
\hline
\textbf{S13}                                                                      & \begin{tabular}[c]{@{}l@{}}E-commerce customer churn prevention using machine learning-based \\ business intelligence strategy \cite{S13-ID202-Gangadhar}\end{tabular}                                                                                          & 0.5        & 0.5        & 0.5        & 1          & 0.5        & 0.5        & \textbf{3.5}     \\ 
\hline
\rowcolor{lightgray} \textbf{S14}                                    & \begin{tabular}[c]{@{}>{\cellcolor{lightgray}}l@{}}Customer Churn Prediction Using Ordinary Artificial Neural Network and Convolutional \\ Neural Network Algorithms: A Comparative Performance Assessment \cite{S14-ID226-Seymen}\end{tabular} & 0.5        & 1          & 0.5        & 0          & 1          & 0.5        & \textbf{3.5}     \\ 
\hline
\textbf{S15}                                                                      & Novel time slicing approach for customer defection models in e-commerce: a case study \cite{S15-ID297-Georgiou}                                                                                                                                               & 1          & 0.5        & 0.5        & 0.5        & 0.5        & 0.5        & \textbf{3.5}   \\ 
\hline
\rowcolor{lightgray} \textbf{S16}                                    & \begin{tabular}[c]{@{}>{\cellcolor{lightgray}}l@{}}Customer Churn in Retail E-Commerce Business: Spatial and Machine Learning \\ Approach \cite{S16-ID325-Matuszelanski} \end{tabular} &                                                                                                                                                1          & 1          & 1          & 0.5        & 1          & 1          & \textbf{5.5}     \\ 
\hline
\textbf{S17}                                                                      & XGBoost-Based E-Commerce Customer Loss Prediction \cite{S17-ID344-Gan}                                                                                                                                                                                    & 1          & 1          & 0.5        & 1          & 1          & 1          & \textbf{5.5}   \\ 
\hline
\textbf{S18}                                                                      & \begin{tabular}[c]{@{}l@{}}An Improved Machine Learning Based Customer Churn Prediction for Insight and \\ Recommendation in E-commerce \cite{S18-ID354-Jahan}\end{tabular}                                                                                 & 0          & 0.5        & 0.5        & 0.5        & 0.5        & 1          & \textbf{3}     \\ 
\hline
\rowcolor{lightgray} \textbf{S19}                                    & Churn Detection Using Machine Learning in the Retail Industry \cite{S19-ID363-Agbemadon}                                                                                                                                                                       & 0          & 1          & 0.5        & 1          & 1          & 1          & \textbf{4.5}   \\ 
\hline
\textbf{S20}                                                                      & \begin{tabular}[c]{@{}l@{}}Performance Evaluation of Various Classification Techniques for Customer Churn \\ Prediction in E-commerce \cite{S20-ID379-Baghla} \end{tabular}                                                                                   & 0.5        & 1          & 1          & 0.5        & 1          & 1          & \textbf{5}     \\ 
\hline
\rowcolor{lightgray} \textbf{S21}                                    & Customer Churn Prediction Using Deep Learning \cite{S21-ID389-Seymen}                                                                                                                                                                                       & 0          & 1          & 0.5        & 0          & 1          & 0.5        & \textbf{3}   \\ 
\hline
\textbf{S22}                                                                      & A PCA-AdaBoost model for E-commerce customer churn prediction \cite{S22-ID394-Wu}                                                                                                                                                                       & 1          & 0.5        & 0.5        & 0.5        & 1          & 1          & \textbf{4.5}   \\ 
\hline
\rowcolor{lightgray} \textbf{S23}                                    & Deep learning for customer churn prediction in e-commerce decision support \cite{S23-ID396-Pondel}                                                                                                                                                          & 0.5        & 1          & 0.5        & 0.5        & 1          & 0.5        & \textbf{4}     \\ 
\hline
\textbf{S24}                                                                      & A Clustering-Prediction Pipeline for Customer Churn Analysis \cite{S24-ID414-Zheng}                                                                                                                                                                        & 0.5        & 0.5        & 1          & 1          & 0.5        & 1          & \textbf{4.5}     \\ 
\hline
\rowcolor{lightgray} \textbf{S25}                                    & A New Time Series Approach in Churn Prediction with Discriminatory Intervals \cite{S25-ID417-Ahmadi}                                                                                                                                                        & 0          & 0.5        & 1          & 0.5        & 1          & 1          & \textbf{4}     \\ 
\hline
\textbf{S26}                                                                      & User Churn Model in E-Commerce Retail \cite{S26-ID442-Fridrich}                                                                                                                                                                                               & 0          & 1          & 1          & 0          & 1          & 1          & \textbf{4}     \\ 
\hline
\rowcolor{lightgray} \textbf{S27}                                    & Customer Churn Prediction in FMCG Sector Using Machine Learning Applications \cite{S27-ID471-Gunesen}                                                                                                                                                        & 0          & 0.5        & 0.5        & 1          & 0.5        & 0.5        & \textbf{3}     \\ 
\hline
\textbf{S28}                                                                      & Predicting Customer Turnover Using Recursive Neural Networks \cite{S28-ID472-Chashmi}                                                                                                                                                                        & 1          & 0.5        & 1          & 0          & 1          & 1          & \textbf{4.5}     \\ 
\hline
\rowcolor{lightgray} \textbf{S29}                                    & User Modeling for Churn Prediction in E-Commerce \cite{S29-ID538-Berger}                                                                                                                                                                                    & 1          & 0.5        & 0.5        & 0.5        & 0.5        & 0.5        & \textbf{3.5}     \\ 
\hline
\textbf{S30}                                                                      & \begin{tabular}[c]{@{}l@{}}Clustering prediction techniques in defining and predicting customers defection: \\ The case of e-commerce context \cite{S30-ID647-Rachid}\end{tabular}                                                                           & 1          & 0.5        & 0.5        & 0.5        & 0.5        & 1          & \textbf{4}     \\ 
\hline
\rowcolor{lightgray} \textbf{S31}                                    & \begin{tabular}[c]{@{}>{\cellcolor{lightgray}}l@{}}Predicting partial customer churn using Markov for discrimination for modeling \\ first purchase sequences \cite{S31-ID717-Migueis}\end{tabular}                                              & 1          & 0.5        & 0.5        & 0.5        & 0.5        & 0.5        & \textbf{3.5}     \\ 
\hline
\textbf{S32}                                                                      & Research on e-commerce customer churning modeling and prediction \cite{S32-ID720-Zhao}                                                                                                                                                                    & 0          & 1          & 0.5        & 0.5        & 0.5        & 0.5        & \textbf{3}   \\ 
\hline
\rowcolor{lightgray} \textbf{S33}                                    & \begin{tabular}[c]{@{}>{\cellcolor{lightgray}}l@{}}Modeling partial customer churn: On the value of first product-category \\ purchase sequences \cite{S33-ID742-Migueis}\end{tabular}                                                           & 1          & 0.5        & 0.5        & 0          & 0.5        & 0.5        & \textbf{3}   \\ 
\hline
\textbf{S34}                                                                      & \begin{tabular}[c]{@{}l@{}}An extended support vector machine forecasting framework for customer \\ churn in e-commerce \cite{S34-ID749-Yu}\end{tabular}                                                                                                 & 1          & 0.5        & 0          & 1          & 0.5        & 0.5        & \textbf{3.5}   \\ 
\hline
\rowcolor{lightgray} \textbf{S35}                                    & \begin{tabular}[c]{@{}>{\cellcolor{lightgray}}l@{}}Customer attrition in retailing: An application of Multivariate Adaptive Regression \\ Splines \cite{S35-ID783-Miguéis}\end{tabular}                                                          & 1          & 0.5        & 0.5        & 0          & 0.5        & 1          & \textbf{3.5}     \\ 
\hline
\textbf{S36}                                                                      & Applying data classification techniques for churn prediction in retailing \cite{S36-ID806-Ching}                                                                                                                                                           & 0.5        & 0.5        & 0.5        & 0          & 0.5        & 1          & \textbf{3}   \\ 
\hline
\rowcolor{lightgray} \textbf{S37}                                    & \begin{tabular}[c]{@{}>{\cellcolor{lightgray}}l@{}}Customer base analysis: Partial defection of behaviourally loyal clients in a \\ non-contractual FMCG retail setting \cite{S37-ID826-Buckinx}\end{tabular}                                    & 1          & 0.5        & 0.5        & 0          & 1          & 1          & \textbf{4}     \\ 
\hline
\textbf{S38}                                                                      & Handling class imbalance in customer churn prediction                                                                                                                                                                                & 1          & 0.5        & 0          & 1          & 1          & 1          & \textbf{4.5}   \\ 
\hline
\rowcolor{lightgray} \textbf{S39}                                    & A Customer Churn Prediction Model for Hartlief Shop \& Bistro, Namibia \cite{S39-ID861-Ashipala}                                                                                                                                                                & 0          & 0.5        & 0.5        & 0.5        & 1          & 1          & \textbf{3.5}     \\ 
\hline
\textbf{S40}                                                                      & Comparative Study of Customer Churn Prediction Based on Data Ensemble Approach \cite{S40-ID1045-Kumar}                                                                                                            & 0          & 0.5        & 0          & 0.5        & 1          & 1          & \textbf{3}     \\ 
\hline
\rowcolor{lightgray} \textbf{S41}                                    & Predicting Customer Churn in Retailing \cite{S41-ID1064-Sweidan}                                                                                                                                                                                              & 0.5        & 0.5        & 0.5        & 0          & 1          & 1          & \textbf{3.5}     \\
\hline
\end{tabular}
}
\end{table}


\begin{figure}[ht]
    \centering
    \includegraphics[width=0.6\textwidth]{images/selection-process.pdf}
    \caption{Studies selection process.}
    \label{fig:selection-process}
\end{figure}